\section{Lecture 11: Model Evaluation and Selection}

\subsection{Motivation}

In previous lectures, we developed models for prediction:
\begin{itemize}
    \item linear and logistic regression,
    \item kernel methods and support vector machines.
\end{itemize}

\medskip

\noindent
However, a fundamental question remains:

\medskip

\noindent
\textit{How do we know whether a model is good?}

\medskip

\noindent
\textbf{Guiding idea.}  
Model evaluation estimates how well a model generalizes to unseen data.

\subsection{Training, Validation, and Testing}

We typically split data into three parts:

\medskip

\noindent
\textbf{Training set.}
\begin{itemize}
    \item Used to fit model parameters
\end{itemize}

\medskip

\noindent
\textbf{Validation set.}
\begin{itemize}
    \item Used to tune hyperparameters
    \item Guides model selection
\end{itemize}

\medskip

\noindent
\textbf{Test set.}
\begin{itemize}
    \item Used only for final evaluation
    \item Provides an unbiased estimate of performance
\end{itemize}

\medskip

\noindent
\textbf{Workflow.}
\begin{enumerate}
    \item Train models on training data
    \item Select best model using validation data
    \item Report final performance on test data
\end{enumerate}

\medskip

\noindent
\textbf{Key principle.}  
The test set must not influence model selection.

\medskip

\noindent
\textbf{Insight.}  
Proper data splitting is essential for reliable evaluation.

\subsection{Cross-Validation}

When data is limited, a single split may be unreliable.

\medskip

\noindent
\textbf{$k$-fold cross-validation.}
\begin{itemize}
    \item Partition data into $k$ subsets (folds)
    \item Train on $k-1$ folds, validate on remaining fold
    \item Repeat for all folds
\end{itemize}

\medskip

\noindent
\textbf{Estimate.}
\[
\text{CV error} = \frac{1}{k} \sum_{j=1}^k \text{error}_j.
\]

\medskip

\noindent
\textbf{Benefits.}
\begin{itemize}
    \item More stable estimate of generalization error
    \item Efficient use of data
\end{itemize}

\medskip

\noindent
\textbf{Tradeoffs.}
\begin{itemize}
    \item Higher computational cost
    \item Choice of $k$ affects bias–variance of estimate
\end{itemize}

\medskip

\noindent
\textbf{Insight.}  
Cross-validation approximates test performance without requiring a large held-out set.

\subsection{Performance Metrics}

The choice of metric depends on the task.

\medskip

\noindent
\textbf{Regression metrics.}
\begin{itemize}
    \item Mean squared error (MSE):
    \[
    \frac{1}{n} \sum (y_i - \hat{y}_i)^2
    \]
    \item Mean absolute error (MAE)
\end{itemize}

\medskip

\noindent
\textbf{Classification metrics.}

\begin{itemize}
    \item Accuracy:
    \[
    \frac{\text{correct predictions}}{\text{total}}
    \]
    \item Precision, recall
    \item F1-score
\end{itemize}

\medskip

\noindent
\textbf{Probabilistic metrics.}
\begin{itemize}
    \item Cross-entropy loss
    \item Log-likelihood
\end{itemize}

\medskip

\noindent
\textbf{Important consideration.}
\begin{itemize}
    \item Metrics should reflect the goal of the task
    \item Different metrics emphasize different aspects of performance
\end{itemize}

\medskip

\noindent
\textbf{Insight.}  
Choosing the right metric is part of defining the learning problem.

\subsection{ROC Curves and AUC}

In many classification problems, especially those involving probabilistic predictions, evaluating performance using a single threshold (e.g., accuracy) can be misleading.

\medskip

\noindent
\textbf{Key idea.}  
Instead of fixing a threshold, evaluate performance across all possible thresholds.

\subsubsection{True and False Positive Rates}

Consider binary classification with labels $y \in \{0,1\}$.

\medskip

\noindent
\textbf{Definitions.}

\begin{itemize}
    \item \textbf{True Positive Rate (TPR)} (Recall):
    \[
    \text{TPR} = \frac{\text{True Positives}}{\text{Actual Positives}}
    \]

    \item \textbf{False Positive Rate (FPR)}:
    \[
    \text{FPR} = \frac{\text{False Positives}}{\text{Actual Negatives}}
    \]
\end{itemize}

\subsubsection{ROC Curve}

The \emph{Receiver Operating Characteristic (ROC)} curve plots:
\[
\text{TPR} \quad \text{vs} \quad \text{FPR}
\]
as the classification threshold varies.

\medskip

\noindent
\textbf{Interpretation.}

\begin{itemize}
    \item Each point corresponds to a different threshold
    \item The curve shows the tradeoff between sensitivity and false alarms
\end{itemize}

\medskip

\noindent
\textbf{Special cases.}

\begin{itemize}
    \item Perfect classifier: reaches $(0,1)$ (top-left corner)
    \item Random classifier: diagonal line $\text{TPR} = \text{FPR}$
\end{itemize}

\medskip

\noindent
\textbf{Insight.}  
ROC curves evaluate ranking quality rather than fixed-threshold decisions.

\subsubsection{Area Under the Curve (AUC)}

The \emph{Area Under the ROC Curve (AUC)} is defined as:
\[
\text{AUC} = \int_0^1 \text{TPR}(\text{FPR}) \, d(\text{FPR}).
\]

\medskip

\noindent
\textbf{Interpretation.}

\begin{itemize}
    \item $\text{AUC} = 1$: perfect classifier
    \item $\text{AUC} = 0.5$: random guessing
\end{itemize}

\medskip

\noindent
\textbf{Probabilistic interpretation.}

\medskip

\noindent
AUC equals the probability that:
\[
\text{score}(x^+) > \text{score}(x^-),
\]
where $x^+$ is a positive example and $x^-$ is a negative example.

\medskip

\noindent
\textbf{Insight.}  
AUC measures how well the model ranks positive examples above negative ones.

\subsubsection{When to Use ROC and AUC}

\textbf{Advantages.}

\begin{itemize}
    \item Threshold-independent evaluation
    \item Useful for imbalanced datasets
    \item Captures ranking performance
\end{itemize}

\medskip

\noindent
\textbf{Limitations.}

\begin{itemize}
    \item Does not reflect specific operating points
    \item May be overly optimistic in highly skewed settings
\end{itemize}

\medskip

\noindent
\textbf{Practical guidance.}

\begin{itemize}
    \item Use ROC/AUC when comparing models broadly
    \item Use precision/recall when class imbalance is critical
\end{itemize}

\subsubsection{Connection to Previous Concepts}

\begin{itemize}
    \item Logistic regression outputs probabilities used for ranking
    \item Cross-entropy optimizes likelihood, not directly AUC
    \item Evaluation metrics define what ``good performance'' means
\end{itemize}

\medskip

\noindent
\textbf{Key insight.}  
Different metrics emphasize different aspects of prediction quality; choosing the right one is essential.

\subsection{Model Comparison}

We often need to choose between multiple models or hyperparameters.

\medskip

\noindent
\textbf{Approach.}
\begin{itemize}
    \item Train multiple models
    \item Evaluate using validation or cross-validation
    \item Select best-performing model
\end{itemize}

\medskip

\noindent
\textbf{Hyperparameters.}
\begin{itemize}
    \item Regularization strength
    \item Kernel parameters
    \item Model complexity
\end{itemize}

\medskip

\noindent
\textbf{Overfitting to validation data.}
\begin{itemize}
    \item Excessive tuning can overfit validation set
    \item Requires careful experimental design
\end{itemize}

\medskip

\noindent
\textbf{Statistical considerations.}
\begin{itemize}
    \item Performance differences may be due to noise
    \item Multiple evaluations increase risk of overfitting
\end{itemize}

\medskip

\noindent
\textbf{Insight.}  
Model selection is itself a learning problem and must be treated carefully.

\subsection{A Unifying Perspective}

Model evaluation connects several aspects of learning:

\begin{itemize}
    \item \textbf{Generalization:} performance on unseen data
    \item \textbf{Optimization:} training minimizes empirical loss
    \item \textbf{Regularization:} controls model complexity
\end{itemize}

\medskip

\noindent
\textbf{Core idea.}  
Evaluation estimates how well learned models approximate the true underlying function.

\medskip

\noindent
\textbf{Key insight.}  
Reliable evaluation requires separating model fitting from model assessment.

\subsection{Outlook}

In this lecture, we studied:

\begin{itemize}
    \item data splitting strategies,
    \item cross-validation,
    \item performance metrics,
    \item model comparison.
\end{itemize}

\medskip

\noindent
These tools are essential for practical machine learning.

\medskip

\noindent
In the next lecture, we return to optimization, focusing on practical aspects such as stochastic gradient descent and modern optimization methods.

\medskip

\noindent
\textbf{Guiding principle.}  
A model is only as good as its performance on unseen data, and evaluation must reflect this.